\section{Prospective Methods and Ethics}

The proposed programme comprises content review, cognitive interviews, a development sample, an independent confirmation sample, and a defined retest subsample. Eligible participants are consenting adults able to understand the authorized language version and report on a current romantic relationship. Final relationship criteria, sample sizes, ceilings, burden estimates, compensation, comparator versions and numeric analysis rules remain to be frozen before the relevant phase. Recruitment through the investigator's own assessment relationships is excluded from the default design.

Experimental items ask about intentional attempts to initiate changes in contact, availability, boundaries and interaction structure over a fixed reference period. Frequency of opportunities, effectiveness, consent and harmfulness are separate variables. No item bank, total score, norm or clinical cutoff is validated by this proposal. Adjacent constructs require authorized instruments and adequate sample-specific measurement quality. Initial criterion evidence uses independently coded neutral vignette responses; daily ESM, partner reporting and dynamic-model estimation belong to M04.

The institution, ethics decision, processing bases, privacy notice, retention, withdrawal mechanism and trained distress contact must be documented before collection. Participation is individual and voluntary, with optional questions, breaks and withdrawal. Default interviews use written notes; no audio, video or physiological collection is proposed. Contact/consent records are separated from pseudonymous research data. No individual attachment, controller or relationship diagnosis is provided. Prepared materials are maintained in \texttt{human\_study/}; their existence does not imply submission or approval.
